\documentclass[10pt]{beamer}

\usetheme{Madrid}
\usecolortheme{whale}

\usepackage[utf8]{inputenc}
\usepackage[T1]{fontenc}
\usepackage[italian]{babel}
\usepackage{tikz}
\usetikzlibrary{arrows.meta, positioning, shapes.geometric, calc, fit}
\usepackage{listings}
\usepackage{xcolor}
\usepackage{booktabs}
\usepackage{graphicx}

\definecolor{primary}{RGB}{38, 66, 139}
\definecolor{codebg}{RGB}{245, 245, 245}

\lstset{
    backgroundcolor=\color{codebg},
    basicstyle=\ttfamily\small,
    breaklines=true,
    frame=single,
    showstringspaces=false
}

\title[Lezione 4]{Programmazione e Tecnologie Web}
\subtitle{Backend Architecture \& API}
\author{Giuseppe Capasso}
\date{2025/2026}

\begin{document}

\begin{frame}
    \titlepage
\end{frame}

\section{Architettura Backend}
\begin{frame}{Il Ruolo del Server}
    Il backend gestisce la logica di business e l'accesso alle risorse.
    \begin{itemize}
        \item \textbf{Controllo}: Gestione del ciclo di vita delle risorse.
        \item \textbf{Accesso}: Interfaccia diretta a file system e calcolo.
        \item \textbf{API-First}: Il backend espone solo dati, non HTML.
    \end{itemize}
    \vspace{0.2cm}
    \textit{Strategia didattica: Persistenza su file JSON per concentrarsi sul protocollo HTTP.}
\end{frame}

\begin{frame}[fragile]{Handler Pattern}
    L'handler associa una rotta a una funzione di gestione.
    \begin{itemize}
        \item Esegue il flusso: Validazione $\rightarrow$ Esecuzione $\rightarrow$ Risposta.
    \end{itemize}
    \vspace{0.5cm}
    \centering
    \begin{tikzpicture}[
        block/.style={rectangle, draw, fill=blue!10, minimum width=2.5cm, minimum height=0.8cm, align=center, font=\small},
        arrow/.style={-latex, thick, primary}
    ]
        \node (v) [block] {Validazione};
        \node (e) [block, right=0.8cm of v] {Esecuzione};
        \node (r) [block, right=0.8cm of e] {Risposta};
        \draw [arrow] (v) -- (e);
        \draw [arrow] (e) -- (r);
    \end{tikzpicture}
\end{frame}

\begin{frame}{Flusso della Richiesta: Router}
    Il Router smista le richieste in base a URL e Metodo HTTP.
    \vspace{0.3cm}
    \centering
    \begin{tikzpicture}[
        request/.style={rectangle, draw, fill=blue!10, minimum width=1.5cm, minimum height=0.6cm, font=\tiny},
        router/.style={diamond, draw, fill=yellow!10, minimum width=2cm, minimum height=1.2cm, align=center, font=\tiny},
        handler/.style={rectangle, draw, fill=primary!10, minimum width=1.5cm, minimum height=0.6cm, font=\tiny},
        file/.style={cylinder, draw, rotate=90, fill=gray!10, minimum width=0.6cm, minimum height=0.8cm, font=\tiny}
    ]
        \node (req) [request] {HTTP Request};
        \node (rot) [router, right=1.2cm of req] {Router};
        \node (h1) [handler, right=1.2cm of rot, yshift=0.8cm] {Handler GET};
        \node (h2) [handler, right=1.2cm of rot, yshift=-0.8cm] {Handler POST};
        \node (f) [file, right=1.2cm of h2] {data.json};

        \draw [-latex] (req) -- (rot);
        \draw [-latex] (rot) -- node[above, scale=0.5] {/tasks} (h1);
        \draw [-latex] (rot) -- node[above, scale=0.5] {/tasks} (h2);
        \draw [-latex] (h2.east) -- (f.north);
    \end{tikzpicture}
\end{frame}

\section{Setup Server}
\begin{frame}[fragile]{Setup: Installazione Flask}
    Preparazione dell'ambiente Python per lo sviluppo backend.
    \begin{itemize}
        \item \textbf{Installazione}:
        \begin{verbatim}
pip install flask
        \end{verbatim}
        \item \textbf{Esecuzione (Debug Mode)}:
        \begin{verbatim}
export FLASK_APP=app.py
flask run --debug
        \end{verbatim}
    \end{itemize}
\end{frame}

\begin{frame}[fragile]{Server Flask (Python)}
    Flask è un micro-framework ideale per prototipare API REST.
    \begin{verbatim}
from flask import Flask, jsonify
app = Flask(__name__)

@app.route('/', methods=['GET'])
def index():
    return jsonify({"msg": "Server Flask attivo!"})
    \end{verbatim}
\end{frame}

\begin{frame}[fragile]{Setup: Progetto Express}
    Inizializzazione di un progetto Node.js e installazione dipendenze.
    \begin{itemize}
        \item \textbf{Inizializzazione}:
        \begin{verbatim}
npm init -y
        \end{verbatim}
        \item \textbf{Installazione}:
        \begin{verbatim}
npm install express
        \end{verbatim}
    \end{itemize}
\end{frame}

\begin{frame}[fragile]{Server Express.js (Node.js)}
    Express è lo standard per gestire middleware e routing in Node.
    \begin{verbatim}
const express = require('express');
const app = express();

app.get('/', (req, res) => {
    res.json({ message: "Server Express attivo!" });
});

app.listen(3000);
    \end{verbatim}
\end{frame}

\section{REST API}
\begin{frame}{Todo App: Definizione Rotte}
    Esempio pratico: Gestione di una lista di attività (Task).
    \begin{itemize}
        \item \textbf{Recupero}: \texttt{GET /tasks}
        \item \textbf{Creazione}: \texttt{POST /tasks}
        \item \textbf{Eliminazione}: \texttt{DELETE /tasks/<id>}
    \end{itemize}
    \vspace{0.3cm}
    \textit{Il backend deve supportare il parsing del corpo della richiesta (JSON).}
\end{frame}

\begin{frame}{Semantica REST}
    Uso corretto dei metodi HTTP:
    \begin{itemize}
        \item \texttt{GET}: Recupero dati (idempotente).
        \item \texttt{POST}: Creazione dati.
        \item \texttt{DELETE}: Eliminazione.
    \end{itemize}
    \vspace{0.3cm}
    \textbf{Codici di Stato}:
    \begin{itemize}
        \item 201 (Created), 400 (Bad Request), 401 (Unauthorized).
    \end{itemize}
\end{frame}

\begin{frame}[fragile]{Implementazione REST (Flask)}
    Gestione della rotta POST per aggiungere task.
    \begin{verbatim}
@app.route('/tasks', methods=['POST'])
def add_task():
    new_task = request.get_json()
    tasks.append(new_task)
    return jsonify({"status": "created"}), 201
    \end{verbatim}
\end{frame}

\begin{frame}[fragile]{Implementazione REST (Express)}
    Gestione della rotta POST con middleware JSON.
    \begin{verbatim}
app.use(express.json());

app.post('/tasks', (req, res) => {
    const newTask = req.body;
    tasks.push(newTask);
    res.status(201).json({ status: "created" });
});
    \end{verbatim}
\end{frame}

\section{Persistenza}
\begin{frame}{Strategie di Persistenza}
    Mantenere i dati tra i riavvii del server.
    \begin{enumerate}
        \item \textbf{Lettura}: Caricamento del file JSON in memoria.
        \item \textbf{Modifica}: Manipolazione della struttura dati (array/dict).
        \item \textbf{Scrittura}: Salvataggio atomico su disco.
    \end{enumerate}
    \vspace{0.2cm}
    \textit{Vantaggio JSON: Leggibile dall'uomo e nativo per il web.}
\end{frame}

\begin{frame}[fragile]{Persistenza JSON}
    L'operazione deve essere atomica: Load $\rightarrow$ Modifica $\rightarrow$ Write.
    \begin{verbatim}
def save_tasks(tasks):
    with open('tasks.json', 'w') as f:
        json.dump(tasks, f, indent=4)
    \end{verbatim}
    \textit{Nota: JSON è ottimo per didattica, ma occhio alla concorrenza.}
\end{frame}

\begin{frame}[fragile]{Persistenza JSON (Express)}
    Utilizzo del modulo \texttt{fs} per la scrittura su file.
    \begin{verbatim}
const fs = require('fs');
const FILE_PATH = './tasks.json';

const saveTasks = (tasks) => {
    const data = JSON.stringify(tasks, null, 4);
    fs.writeFileSync(FILE_PATH, data);
};
    \end{verbatim}
\end{frame}

\section{Sicurezza}
\begin{frame}[fragile]{Middleware di Sicurezza}
    Filtrare le richieste prima dell'esecuzione logica.
    \vspace{0.5cm}
    \centering
    \begin{tikzpicture}[
        node/.style={rectangle, draw, fill=red!10, minimum width=2cm, minimum height=0.8cm, align=center, font=\small},
        arrow/.style={-latex, thick, primary}
    ]
        \node (auth) [node] {Autenticazione};
        \node (val) [node, right=0.8cm of auth] {Validazione};
        \draw [arrow] (auth) -- (val);
    \end{tikzpicture}
    \begin{itemize}
        \item \textbf{Basic Auth}: Header codificato Base64.
        \item \textbf{Sanitizzazione}: Pulizia input obbligatoria.
    \end{itemize}
\end{frame}

\begin{frame}{Validazione e Sanitizzazione}
    La regola aurea: \textbf{Mai fidarsi dell'input del client}.
    \begin{itemize}
        \item \textbf{Validazione}: Verifica che i dati siano nel formato corretto (es. l'ID deve essere un numero).
        \item \textbf{Sanitizzazione}: Rimozione di tag script o caratteri pericolosi per prevenire iniezioni.
    \end{itemize}
    \textit{Questi processi avvengono prima di processare la business logic.}
\end{frame}

\begin{frame}[fragile]{Basic Auth (Flask)}
    Proteggere un endpoint con credenziali.
    \begin{verbatim}
@app.route('/secure')
def secure_endpoint():
    auth = request.authorization
    if not auth or auth.password != 'secret':
        return Response('Accesso negato', 401, 
           {'WWW-Authenticate': 'Basic realm="Login"'})
    return "Dati protetti"
    \end{verbatim}
\end{frame}

\begin{frame}[fragile]{Basic Auth (Express)}
    Utilizzo di middleware per la protezione globale o locale.
    \begin{verbatim}
const basicAuth = require('express-basic-auth');

app.use(basicAuth({
    users: { 'admin': 'secret' },
    challenge: true
}));
    \end{verbatim}
\end{frame}

\section{Frontend}
\begin{frame}[fragile]{Fetch API Standard}
    Eseguire richieste asincrone senza autenticazione.
    \begin{verbatim}
fetch('/tasks')
    .then(response => response.json())
    .then(data => console.log(data))
    .catch(error => console.error('Errore:', error));
    \end{verbatim}
\end{frame}

\begin{frame}[fragile]{Integrazione: Fetch API}
    Consumare le API asincronamente.
    \begin{verbatim}
// Autenticazione dinamica via Prompt
const creds = btoa(`${prompt('User:')}:${prompt('Pass:')}`);
fetch('/secure', {
    headers: { 'Authorization': 'Basic ' + creds }
})
.then(r => r.json());
    \end{verbatim}
\end{frame}

\begin{frame}{Conclusioni ed Esercizi}
    \textbf{Obiettivi Raggiunti}:
    \begin{itemize}
        \item Creazione di un server Web (Flask/Express).
        \item Progettazione di API RESTful.
        \item Persistenza dati su JSON.
    \end{itemize}
    \vspace{0.3cm}
    \textbf{Esercizio}: Implementare la cancellazione di un task (\texttt{DELETE /tasks/<id>}) aggiornando il file JSON.
\end{frame}

\end{document}
